%=========================================================================
% appendix.tex
% Supplementary analyses for "Social Vulnerability and Progression to
% Vision-Threatening Diabetic Retinopathy"
%
% Each test is presented in four parts, as agreed:
%   (1) What is tested, conceptually
%   (2) How we tested it
%   (3) What we found
%   (4) Interpretation
%=========================================================================
\documentclass[11pt]{article}

\usepackage[margin=1in]{geometry}
\usepackage{setspace}
\usepackage{booktabs}
\usepackage{amsmath}
\usepackage{graphicx}
\usepackage{tikz}
\usetikzlibrary{arrows.meta, positioning, calc}
\usepackage[hidelinks]{hyperref}
\usepackage{titlesec}
\usepackage{tocloft}
% \thesection is "Test~N", which is far wider than the default TOC number box
% and caused the number to overlap the first word of each entry.
\setlength{\cftsecnumwidth}{5.2em}
\setlength{\cftsecindent}{0em}
\renewcommand{\cftsecfont}{\normalfont}
\renewcommand{\cftsecpagefont}{\normalfont}

\onehalfspacing
\renewcommand{\thesection}{Test~\arabic{section}}
\titleformat{\section}{\large\bfseries}{\thesection.}{0.6em}{}
\newcommand{\tpart}[1]{\medskip\noindent\textbf{#1}\quad}

\title{\textbf{Supplementary Appendix}\\[6pt]
\large Companion analyses for:\\[2pt]
\emph{Social Vulnerability and Progression to Vision-Threatening Diabetic
Retinopathy}\\[2pt]
\emph{Proximity to Academic Ophthalmology and the Severity of Diabetic
Retinopathy}\\[4pt]
\normalsize A county-level analysis of 3,119 United States counties, 2021}
\date{}

\begin{document}
\maketitle

\noindent This appendix documents every hypothesis we tested across both
companion manuscripts, including those that did not support either. Each is
presented as: what is tested conceptually, how we tested it, what we found, and
how we interpret it.

\medskip
\noindent Tests 1, 3, 5, 6, 9, 10, 11 and 12 support the disparities manuscript
(social vulnerability and progression). Tests 2, 13 and 14 are the evidence base
for the companion manuscript on academic ophthalmology proximity. Tests 4, 7 and
8 are diagnostics that constrain the interpretation of both. Every analysis
reported in either manuscript appears here.

\medskip
\noindent Unless stated otherwise, the outcome is the vision-threatening share
of diagnosed diabetic retinopathy (VTDR cases divided by all-stage DR cases
among adults with diabetes), models are quasibinomial with DR cases as trials,
and continuous predictors are standardized so coefficients are per standard
deviation on the log-odds scale.

\tableofcontents
\clearpage

%=========================================================================
\section{Post-stratification artifact in county surveillance estimates}
%=========================================================================

\tpart{(1) What is tested.}
County-level estimates from the Vision and Eye Health Surveillance System
(VEHSS) are not direct measurements. They are constructed by computing national
prevalence by age, sex, and race, applying those national rates to each county's
demographic composition, and then adjusting by a county random effect estimated
from Medicare and Medicaid claims. If that is so, then a county's racial and age
composition moves its published estimate \emph{by construction}, and any
regression of these estimates on a county-level social index that contains a
demographic component will partly recover the post-stratification machinery
rather than a local epidemiologic signal. The Social Vulnerability Index
contains exactly such a component (Theme 3, racial and ethnic minority status).
We test how much of the apparent social gradient is this artifact.

\tpart{(2) How we tested it.}
We obtained VEHSS county estimates separately for each race/ethnicity by age
stratum (non-Hispanic White, non-Hispanic Black, and Hispanic adults, each aged
40--64 and 65--84). Within a single stratum, demographic composition is constant
by construction, so the only county-varying component is the county random
effect, which is the genuinely local signal. We estimated the SVI coefficient
(a) on the composition-varying ``All races, all ages'' file, as a conventional
analysis would, and (b) within each fixed stratum, and compared them.

\tpart{(3) What we found.}
The composition-varying estimate was $\beta_{\mathrm{SVI}} = 0.146$. The
within-stratum estimates were roughly one fifth as large and remarkably
consistent across strata:

\begin{center}
\begin{tabular}{lrr}
\toprule
Specification & $\beta_{\mathrm{SVI}}$ & \\
\midrule
All races, all ages (composition varies)   & 0.146 & \\
Non-Hispanic White 65--84                  & 0.033 & \\
Non-Hispanic Black 65--84                  & 0.033 & \\
Hispanic 65--84                            & 0.035 & \\
Non-Hispanic White 40--64                  & 0.035 & \\
\midrule
Pooled, stratum fixed effects              & 0.035 & \\
\bottomrule
\end{tabular}
\end{center}

Approximately 76\% of the apparent gradient is attributable to demographic
composition, and decomposition shows this is almost entirely \emph{racial}
composition: holding age fixed alone leaves the coefficient at 0.150, while
holding race fixed alone gives 0.034, essentially the fully stratified value. Overdispersion also fell sharply, from a dispersion parameter of
about 15 in the unstratified model to 1.93 in the within-stratum model,
indicating that much of the excess county-level variance was itself
compositional.

\tpart{(4) Interpretation.}
This is a general caution for the use of VEHSS county estimates, not a finding
about diabetic retinopathy. Studies relating these estimates to composite social
indices, area deprivation measures, or racial composition will substantially
overstate local gradients unless demographic composition is held fixed. Notably,
the artifact is specific to demographic predictors: VEHSS does not use poverty
or other socioeconomic variables as model inputs, so adjustment for the
non-demographic content of SVI is not circular. Within-stratum estimation is a
cheap and effective safeguard and we recommend it as routine.

%=========================================================================
\section{Is proximity to academic ophthalmology uniquely protective?}
%=========================================================================

\tpart{(1) What is tested.}
The motivating hypothesis of the prior version of this work was that counties
distant from academic ophthalmology programs face barriers to subspecialty
retinal care, and therefore worse diabetic eye outcomes. The identification
problem is that academic programs are sited where ophthalmologists already
practice. Raw distance to an academic program therefore conflates proximity to
academic tertiary capacity with proximity to eye care in general. We test
whether anything is attributable to academic centers specifically.

\tpart{(2) How we tested it.}
Four tests of increasing stringency.

\emph{2a. Conditional effect.} We computed a second exposure, drive time from
the population-weighted centroid to the nearest county containing at least one
patient-care ophthalmologist, and entered both exposures jointly. The academic
coefficient is then the association holding general eye care proximity fixed.

\emph{2b. Discordant counties.} A $2\times2$ comparison of counties with versus
without a local ophthalmologist, crossed with above versus below median drive
time to an academic program. This is the comparison a reviewer of the prior
submission requested directly.

\emph{2c. Effect modification.} An interaction between academic drive time and
having no local ophthalmologist. If academic centers supply something
distinctive, their reach should matter most where no local alternative exists.

\emph{2d. Capacity rather than proximity.} A gravity-style access score summing
program capacity (ophthalmology residents) discounted by a Gaussian decay in
drive time with a 120-minute bandwidth, replacing nearest-program distance.

\emph{2e. Program type.} Distance to the nearest medical-school-based program
and to the nearest hospital or health-system program, entered jointly in the
fully adjusted model. A program was designated medical-school-based when its
sponsoring institution is a university or degree-granting medical school (105
programs), and hospital or health-system based otherwise (18, including four
military programs). If the mechanism is subspecialty depth rather than the
presence of trainees, medical-school proximity should dominate.

\tpart{(3) What we found.}

\emph{2a.} Academic drive time: $-0.024$ alone, $-0.023$ conditional on supply
distance. Distance to any ophthalmologist: $+0.029$ alone, $+0.014$ conditional
($p = 0.06$). Both exposures are oriented so that higher values mean worse
access. The two correlate only $0.41$. The signs are opposite:
distance to \emph{any} ophthalmologist behaves as an access hypothesis predicts,
while distance to an \emph{academic} program does the reverse.

\emph{2b.} The $2\times2$:
\begin{center}
\begin{tabular}{lrrr}
\toprule
 & $n$ & VTDR share & Median academic drive (min) \\
\midrule
Local ophthalmologist, far from academic  & 396   & 18.6\% & 183 \\
Local ophthalmologist, near academic      & 773   & 19.6\% & 61  \\
No local ophthalmologist, far from academic & 1,163 & 16.9\% & 195 \\
No local ophthalmologist, near academic   & 787   & 17.2\% & 87  \\
\bottomrule
\end{tabular}
\end{center}
Counties with ample local ophthalmology but far from an academic center do not
show worse case mix than those near one (18.6\% vs 19.6\%). The larger
difference is horizontal: having any local ophthalmologist is associated with a
roughly two point \emph{higher} vision-threatening share.

\emph{2c.} Interaction $\beta = 0.008$, $p = 0.70$. A clean null.

\emph{2d.} Capacity-weighted access $\beta = +0.018$, $p<0.001$ in the
within-stratum model. Because a higher access score denotes \emph{better}
access, the positive sign carries the same meaning as the negative drive-time
coefficient: better modeled access, more vision-threatening disease.

In the spatial model, academic drive time retains a
small negative posterior mean ($-0.010$, 95\% CrI $-0.019$ to $-0.001$): still
signed opposite to the access hypothesis.

\tpart{(4) Interpretation.}
No specification supports the hypothesis, including the two most favorable to it
(capacity weighting and the no-local-alternative subgroup). The consistent
pattern is that proximity to specialist eye care, and especially to
medical-school programs where subspecialty retinal diagnosis concentrates,
predicts \emph{more} diagnosed vision-threatening disease.

\medskip
\noindent\emph{A subdivision we considered and abandoned.} An earlier draft
split the roster into ``university'' and ``community hospital'' programs and
reported that only the former tracked disease severity. We have removed that
analysis. All 123 programs are ACGME-accredited ophthalmology residencies:
every one of them is academic ophthalmology, and institutions such as Cleveland
Clinic, Henry Ford, Geisinger and Cook County are major academic training
centers that no defensible rule would call ``community.'' The apparent contrast
was also unstable, moving from a clean null for the non-university group under
a mechanical name regex to a clearly non-zero coefficient under every
classification we regarded as more accurate. Because the categories do not
correspond to a real distinction in academic ophthalmology, we report the
roster as a single exposure.

\medskip
\noindent Two cautions on the internal evidence. Tests 2a and 2d are not independent:
drive time and the gravity access score encode the same geography with opposite
orientation and correlate at $0.61$, so their agreement is arithmetic rather
than corroborative.  This is difficult to read as academic
centers improving outcomes and easier to read as diagnostic intensity: where the
subspecialists are, vision-threatening disease is more completely identified and
coded. The authors of the underlying surveillance estimates anticipated exactly
this, cautioning that claims diagnoses are influenced by access to
ophthalmologists and that their model may mistake such factors for true
variation in prevalence. We regard test 2e as the closest thing in these data to
a direct measurement of that phenomenon, while noting that it is an
interpretation we have not independently isolated.

%=========================================================================
\section{Does the social gradient operate through insurance coverage?}
%=========================================================================

\tpart{(1) What is tested.}
If the social gradient in disease progression is produced by uninsurance and its
downstream effects on screening and follow-up, then it should be substantially
attenuated in a population with near-universal coverage. Medicare eligibility at
age 65 provides that contrast. This also converts a reviewer objection to the
prior submission into a design: the reviewer argued that restricting to older
adults ``enriches the population more likely to have health insurance,'' which
is precisely what makes the comparison informative.

\tpart{(2) How we tested it.}
We estimated the SVI coefficient separately for adults aged 40--64 and 65--84
with race held fixed by fixed effects, using identical covariate sets, and
tested the difference with a two-sample $z$ statistic.

\tpart{(3) What we found.}
\begin{center}
\begin{tabular}{lrrr}
\toprule
Age band & Mean VTDR share & $\beta_{\mathrm{SVI}}$ & SE \\
\midrule
40--64 years & 20.3\% & 0.037 & 0.0022 \\
65--84 years & 18.0\% & 0.033 & 0.0023 \\
\bottomrule
\end{tabular}
\end{center}
Difference $0.004$; $z = 1.35$; $p = 0.18$. The same pattern holds within each
race stratum (main text, Table 2): 0.034 versus 0.032 among non-Hispanic White
adults, 0.037 versus 0.035 among non-Hispanic Black adults, 0.046 versus 0.037
among Hispanic adults.

\medskip
\noindent\emph{Correction of frame.} An earlier draft estimated this test in the
composition-varying ``All races'' frame, giving 0.150 versus 0.142 (difference
0.008, $p = 0.16$). Those coefficients are inflated by the post-stratification
artifact documented in Test~1 and were inconsistent with this study's primary
design. The values above hold race fixed. The conclusion is unchanged: the
gradient is not narrower where coverage is near-universal.

\tpart{(4) Interpretation.}
Near-universal insurance coverage does not narrow the social gradient in how
advanced diabetic eye disease is. Whatever produces this gradient is therefore
not principally a coverage mechanism. Remaining candidates include
transportation and appointment access, competing demands on patient time,
differential dilated examination rates within insured populations, and
continuity of the primary care relationship through which retinal screening is
ordered. We note the Hispanic 40--64 versus 65--84 contrast (0.046 vs 0.037) is
the largest of the three and is the one comparison where a coverage effect
cannot be excluded.

%=========================================================================
\section{Which dimension of social vulnerability?}
%=========================================================================

\tpart{(1) What is tested.}
SVI decomposes into four themes: socioeconomic status, household
characteristics, racial and ethnic minority status, and housing type and
transportation. Theme 4 contains a no-vehicle-available component and is the
closest thing in the index to a transport access measure. If progression tracked
Theme 4 more than Theme 1, that would be an access finding the geographic
distance measures failed to capture, and it would be directly actionable.

\tpart{(2) How we tested it.}
Each theme was entered separately, and then all four jointly, with the standard
covariate set.

\tpart{(3) What we found.}
Entered separately, the racial and ethnic minority theme dominated
($\beta = 0.246$) against socioeconomic status ($0.135$), housing and
transportation ($0.096$), and household characteristics ($0.080$). Entered
jointly, minority status was $0.225$ while housing and transportation turned
slightly negative ($-0.026$).

\tpart{(4) Interpretation.}
\textbf{We do not report this as a substantive finding, and it should not be
cited as one.} These models were estimated on composition-varying estimates, and
Test~1 establishes that county racial composition moves the VEHSS outcome
mechanically. The apparent dominance of the minority status theme is therefore
largely the post-stratification artifact reappearing, since Theme 3 is
essentially a measure of county racial composition. A properly stratified
decomposition would require SVI themes to vary within race strata, which they do
only weakly. We report the result here solely to document that we ran it and to
explain why it cannot be interpreted. What can be said is the negative point:
there is no evidence that the transportation dimension drives the gradient,
which is consistent with the null geographic access results in Test~2.

%=========================================================================
\section{Does the gradient survive adjustment for detection?}
%=========================================================================

\tpart{(1) What is tested.}
Because the outcome is built from diagnosed disease in claims, the severity
share could rise wherever mild disease goes undetected, independent of true
progression. If the social gradient were simply a detection gradient, it should
attenuate substantially once detection proxies are included.

\tpart{(2) How we tested it.}
We compared the SVI coefficient with no adjustment against the coefficient after
adding ophthalmologist and optometrist density, an indicator for having no local
ophthalmologist, drive time to the nearest ophthalmologist, drive time to the
nearest academic program, population density, and rurality.

\tpart{(3) What we found.}
SVI alone $0.163$; with detection proxies $0.145$, an attenuation of
approximately 11\%. In the within-stratum specification the coefficient is
$0.035$ with the same proxies included.

\tpart{(4) Interpretation.}
The gradient is not primarily a detection artifact under the proxies available
to us, but this test is weak evidence. Provider supply is an imperfect proxy for
screening intensity, and the VEHSS county random effects were themselves already
adjusted for ophthalmologists per capita, so our adjustment is partly a second
pass at a factor already partially removed. Residual detection bias remains the
principal threat to the main finding, and its likely direction is to overstate
it.

%=========================================================================
\section{Spatial dependence}
%=========================================================================

\tpart{(1) What is tested.}
County health outcomes are spatially clustered. If unmeasured regional factors
such as state Medicaid policy, regional diabetes management norms, or health
system catchments drive both social vulnerability and disease progression, an
aspatial model will overstate precision and may bias coefficients.

\tpart{(2) How we tested it.}
Moran's $I$ was computed on Pearson residuals using a queen contiguity weights
matrix over 2021 TIGER/Line county boundaries. Three island counties had no
queen neighbor; each was linked to its nearest county by centroid distance so
the adjacency graph was connected, a decision recorded here rather than left
implicit. We then refitted the largest single stratum (non-Hispanic White,
65--84, 3,106 counties) with a Leroux conditional autoregressive random effect
using CARBayes, with binomial likelihood, VTDR cases as successes and DR cases
as trials, 20,000 burn-in and 140,000 iterations thinned by 20.

\tpart{(3) What we found.}
Moran's $I = 0.237$ ($p < 0.001$), confirming substantial residual clustering.
Estimated spatial dependence was strong ($\rho = 0.82$, 95\% CrI 0.73 to 0.91).
The social vulnerability coefficient attenuated from 0.032 to 0.021 but its
credible interval excluded zero (0.012 to 0.029). Academic drive time was
$-0.010$ (95\% CrI $-0.019$ to $-0.001$), marginally excluding zero but signed
opposite to the access hypothesis. Drive time to the nearest ophthalmologist,
ophthalmologist density, and population density all had credible intervals
spanning zero.

\tpart{(4) Interpretation.}
Roughly a third of the within-stratum social gradient is attributable to broad
spatial structure, which is consistent with regional policy and health system
factors contributing to it. What remains is a genuinely local association that
is not explained by demographic composition, provider supply, rurality, or
spatial position. The contrast between the two exposures under the same spatial
model is the cleanest single statement of the paper's result: social
vulnerability predicts more advanced disease, whereas distance from academic
ophthalmology predicts marginally \emph{less} advanced disease, which is not the
access effect the prior framing proposed.

%=========================================================================
\section{Sensitivity analyses}
%=========================================================================

\tpart{(1) What is tested.}
Whether the main result depends on specific analytic choices: the age handling
criticized in review, the choice of severity share over disease level, and the
possibility that any ratio constructed from VEHSS estimates would show these
patterns.

\tpart{(2) How we tested it.}
\emph{6a. Age standardization.} We computed an age-standardized severity share
weighting stratum-specific shares by the national distribution of DR cases
across VEHSS age strata, and compared it with the crude all-ages share.

\emph{6b. Level versus severity.} We refitted the identical specifications with
diagnosed DR prevalence among adults with diabetes as the outcome, to compare
against the severity share.

\emph{6c. Parallel construction on a different disease.} Glaucoma carries an
analogous severity split in VEHSS (vision-affecting versus all glaucoma). We
built the identical ratio and refitted, asking whether these patterns are
specific to diabetic retinopathy or general to VEHSS-derived severity ratios.

\tpart{(3) What we found.}
\emph{6a.} Crude and age-standardized shares correlate at $r = 0.992$. Results
are unchanged.

\emph{6b.} With the proper diabetes denominator and full covariate adjustment,
drive time was not associated with DR level ($\beta = -0.005$, $p = 0.059$), and
SVI was associated only weakly ($\beta = 0.025$) compared with its association
with severity ($\beta = 0.145$ unstratified). Disease occurrence and disease
progression behave differently, and the severity share is essentially
uncorrelated with the level (Spearman $\rho = -0.08$).

\emph{6c.} The glaucoma severity share has roughly one tenth the relative
variation of the VTDR share (coefficient of variation 0.024 vs 0.231) and
correlates with it at 0.55.

\tpart{(4) Interpretation.}
\emph{6a} retires the age-cutoff objection: the result does not depend on where
an age threshold is placed, because no threshold is used and standardization
does not move it.

\emph{6b} is the substantive justification for the paper's reframing. The social
gradient in progression is several times the gradient in occurrence among people
with diabetes, and the two outcomes are nearly orthogonal, so this is not the
established prevalence finding restated.

\emph{6c} must be reported carefully. We originally proposed glaucoma as a
negative control. It cannot serve that purpose: its severity share barely varies
and it correlates substantially with the DR severity share, so a null there
would reflect low variance rather than specificity. The shared variation is
itself suggestive of a common county-level factor, plausibly care intensity,
which is consistent with the detection interpretation in Test~2 but does not
establish it. A stronger design would use a placebo \emph{exposure}, such as
distance to accredited programs in a specialty unrelated to eye care, which we
did not have data to construct.

%=========================================================================
\section{Detection stress test using Medicare claims volume}
%=========================================================================

\tpart{(1) What is tested.} The principal threat to the main finding is that
mild diabetic retinopathy goes undiagnosed more often in disadvantaged counties,
shrinking the denominator and inflating the vision-threatening share without any
true difference in progression. Provider density (Test~5) is a weak proxy for
this, because having ophthalmologists nearby is not the same as diagnostic
activity occurring. We wanted a measure of how much diagnostic work actually
happens in a county, drawn from claims rather than from workforce counts.

\tpart{(2) How we tested it.} From the CMS Medicare Geographic Variation Public
Use File (county, 2021) we took imaging events, evaluation and management
events, and diagnostic test events per 1,000 Original Medicare beneficiaries.
These were added to the within-stratum specification, singly and jointly, on a
common complete-case sample (13,035 cells).

\tpart{(3) What we found.}
\begin{center}
\begin{tabular}{lrrr}
\toprule
Specification & $\beta_{\mathrm{SVI}}$ & SE & $p$ \\
\midrule
Base (same sample)                & 0.0353 & 0.0016 & $<0.001$ \\
$+$ Medicare imaging per 1,000    & 0.0371 & 0.0016 & $<0.001$ \\
$+$ Medicare E\&M per 1,000       & 0.0375 & 0.0016 & $<0.001$ \\
$+$ all three                     & 0.0354 & 0.0016 & $<0.001$ \\
\bottomrule
\end{tabular}
\end{center}
The coefficients on the detection proxies themselves were imaging $-0.028$
($p<0.001$), E\&M $+0.010$ ($p<0.001$), and tests $+0.007$ ($p = 0.002$).

\tpart{(4) Interpretation.} The social gradient is completely unmoved. What
makes this informative rather than merely null is that the detection pathway is
plainly visible in the same model: counties with more Medicare imaging have a
\emph{lower} vision-threatening share, exactly as differential detection
predicts, since more diagnostic activity finds more mild disease. So the
mechanism exists, is measurable here, and does not account for the social
gradient. This is the strongest evidence we can offer against the detection
explanation. Its limit is that these are all-cause diagnostic volumes, not
retinal screening specifically; a claims measure of dilated eye examination
rates by county would be a better instrument and we could not obtain one at
county level.

%=========================================================================
\section{Preventive-screening engagement as a mechanism}
%=========================================================================

\tpart{(1) What is tested.} If the gradient is not coverage and not detection,
the remaining candidate on our causal diagram is screening cadence and retention
in care. We wanted a measure of whether a county's population actually completes
recommended preventive screening, with no ophthalmic content, so that any
attenuation could not be attributed to eye care access.

\tpart{(2) How we tested it.} County mammography screening completion (County
Health Rankings 2021, derived from Medicare claims) and influenza vaccination
were added to the within-stratum specification.

\tpart{(3) What we found.} The gradient fell from 0.0353 to 0.0206, an
attenuation of approximately 42\% on the log-odds scale and 29\% on the
standardized scale. Mammography completion was itself strongly associated with
the outcome ($\beta = -0.039$, $p<0.001$), and correlated $-0.42$ with social
vulnerability. Influenza vaccination added little ($\beta = 0.005$,
$p = 0.04$).

\tpart{(4) Interpretation.} Roughly two fifths of the social gradient in
diabetic retinopathy progression is shared with general preventive-screening
follow-through. Read alongside Test~8, where diagnostic volume explained none of
it, the pattern is specific: what matters is not how much medical looking occurs
in a county but whether its population completes recommended preventive care.
That points at screening cadence rather than at ophthalmology capacity, and it
is consistent with the null geographic access results in Test~2.

Two cautions. Mammography is a female-only screen applied here to a both-sex
outcome, so it functions as a county-level index of screening culture, not an
individual-level mediator. And because screening engagement is plausibly on the
causal path rather than a confounder, this is a decomposition, not an adjustment;
we do not claim a formal mediation estimate.

%=========================================================================
\section{Direct coverage measures}
%=========================================================================

\tpart{(1) What is tested.} Test~3 inferred the role of insurance indirectly,
from the age contrast at Medicare eligibility. Here we test coverage directly.

\tpart{(2) How we tested it.} The county uninsured rate (County Health Rankings)
and the share of Medicare beneficiaries dually eligible for Medicaid (CMS) were
added to the within-stratum specification. Dual eligibility is informative
because within the 65 and over population everyone already has Medicare, so it
isolates poverty among the fully insured.

\tpart{(3) What we found.} The gradient moved from 0.0353 to 0.0343 with both
included, and neither covariate was independently associated with the outcome
(uninsured $p = 0.14$; dual-eligible $p = 0.44$). Restricting to the 65--84 band
and using dual eligibility as the exposure, however, gave $\beta = 0.006$
($p<0.001$).

\tpart{(4) Interpretation.} Direct coverage measures confirm what the age
contrast implied: insurance status does not explain the social gradient in
progression. The 65--84 dual-eligibility result is the useful positive: among a
population that is universally insured, poverty still predicts more advanced
disease. Coverage is necessary but plainly not sufficient.

%=========================================================================
\section{Medicaid expansion, and a placebo that refutes it}
%=========================================================================

\tpart{(1) What is tested.} Medicaid expansion changed coverage for low-income
adults under 65 and left the 65 and over population untouched. If coverage drove
the social gradient, the gradient among adults 40--64 should be smaller in
expansion states. The 65--84 band is a built-in placebo: expansion cannot have
affected it, so any contrast appearing equally there is confounding.

\tpart{(2) How we tested it.} Counties were classified by 2021 state expansion
status (Kaiser Family Foundation). Missouri and Oklahoma implemented mid-2021
and were excluded rather than forced into either group. We estimated the SVI
gradient in each age-by-expansion cell and fitted the formal triple interaction.

\tpart{(3) What we found.}
\begin{center}
\begin{tabular}{llrr}
\toprule
Age band & Expansion status & Cells & $\beta_{\mathrm{SVI}}$ \\
\midrule
40--64 & Non-expansion & 2,736 & 0.0163 \\
40--64 & Expansion     & 3,822 & 0.0461 \\
65--84 & Non-expansion & 2,436 & 0.0122 \\
65--84 & Expansion     & 3,307 & 0.0421 \\
\bottomrule
\end{tabular}
\end{center}
Triple interaction SVI $\times$ expansion $\times$ age band: $\beta = 0.0023$,
$p = 0.73$.

\tpart{(4) Interpretation.} Read carelessly, the first two rows say Medicaid
expansion triples the social gradient, which would be a striking and
counterintuitive policy finding. The placebo refutes it. The same contrast
appears, essentially unchanged, in the Medicare-age population that expansion
could not have touched, and the triple interaction is null. The difference is
therefore between expansion and non-expansion \emph{states}, which differ in
region, urbanisation, health system structure, and the distribution of social
vulnerability itself, and not a consequence of expansion. We report this test
principally as a demonstration that the placebo design works, and as a caution
against the two-cell comparison that a reader might otherwise draw from these
data.

%=========================================================================
\section{Convergent validity against an independent care-quality outcome}
%=========================================================================

\tpart{(1) What is tested.} Whether the social gradient we report sits within a
coherent pattern of county care-quality gradients, or is anomalous and therefore
more likely to be a measurement artifact of VEHSS.

\tpart{(2) How we tested it.} We regressed two outcomes drawn entirely outside
VEHSS on social vulnerability with the same covariate set: preventable hospital
stays (ambulatory-care-sensitive admissions, Medicare claims) and mammography
screening completion. All outcomes were placed in standard deviation units so
magnitudes are comparable with the vision-threatening share.

\tpart{(3) What we found.} Per standard deviation of social vulnerability:
preventable hospital stays $+0.279$ SD ($p<0.001$); mammography completion
$-0.465$ SD ($p<0.001$); vision-threatening share $+0.100$ SD ($p<0.001$). The
Spearman correlation between preventable hospital stays and the
vision-threatening share was $0.134$.

\tpart{(4) Interpretation.} Social vulnerability predicts ambulatory care
failure in a different organ system and a different data system, which supports
the care-quality reading of our result. It also puts the finding in proportion:
the diabetic retinopathy progression gradient is roughly a third the magnitude
of the preventable-hospitalisation gradient and about a fifth of the
mammography gradient. It is real and consistent, but it is a modest gradient,
and the low correlation between the two outcomes indicates they are related
rather than interchangeable.

%=========================================================================
\section{Dose-response in drive time to academic ophthalmology}
%=========================================================================

\tpart{(1) What is tested.} Test~2 asked whether any academic-proximity effect
is specific to academic centers. This test asks the prior and simpler question:
is there a dose-response at all? A regression coefficient can be driven by a
handful of extreme counties or by a non-monotonic pattern, so we examine the
raw gradient before modelling it.

\tpart{(2) How we tested it.} Counties were binned by drive time to the nearest
residency program accredited by 2021, and the mean county vision-threatening
share computed within each bin, with confidence intervals reflecting
between-county variation.

\tpart{(3) What we found.}
\begin{center}
\begin{tabular}{lrrr}
\toprule
Drive time (minutes) & Counties & Mean VTDR share (\%) & SE \\
\midrule
0--30    & 164 & 18.89 & 0.27 \\
31--60   & 356 & 17.44 & 0.20 \\
61--90   & 464 & 16.94 & 0.18 \\
91--120  & 521 & 16.67 & 0.17 \\
121--180 & 730 & 16.76 & 0.15 \\
181--240 & 370 & 16.19 & 0.21 \\
$>$240   & 514 & 16.12 & 0.17 \\
\bottomrule
\end{tabular}
\end{center}

\tpart{(4) Interpretation.} There is a dose-response, and it runs the wrong way.
The severity of diagnosed disease falls monotonically as distance from an
academic program increases, by roughly 2.8 percentage points from the nearest to
the farthest bin. The access hypothesis predicts the opposite slope. The pattern
is smooth rather than driven by outliers, which rules out the possibility that
the negative regression coefficient is an artifact of a few extreme counties. It
is, however, exactly what differential detection predicts: where subspecialty
capacity is concentrated, more mild disease is found and the denominator grows.

%=========================================================================
\section{Stability of the academic association under adjustment}
%=========================================================================

\tpart{(1) What is tested.} A confounded association should move toward zero as
the confounders are adjusted away. A structural feature of the measurement
should not. We therefore track the academic-distance coefficient itself, rather
than the social vulnerability coefficient, across the same progressive
adjustment sequence used in Tests~8 and~9.

\tpart{(2) How we tested it.} The coefficient on log drive time to the nearest
residency program was estimated on a single common complete-case sample (13,008
cells) under five nested specifications, adding general eye care proximity and
geography, then social vulnerability, then Medicare diagnostic volume, then
preventive-screening completion.

\tpart{(3) What we found.}
\begin{center}
\begin{tabular}{lrrr}
\toprule
Specification & $\beta_{\mathrm{academic}}$ & SE & $p$ \\
\midrule
Unadjusted                              & $-0.0121$ & 0.0007 & $<0.001$ \\
$+$ eye care supply and geography       & $-0.0217$ & 0.0010 & $<0.001$ \\
$+$ social vulnerability                & $-0.0174$ & 0.0010 & $<0.001$ \\
$+$ Medicare diagnostic volume          & $-0.0166$ & 0.0010 & $<0.001$ \\
$+$ preventive screening completion     & $-0.0133$ & 0.0010 & $<0.001$ \\
\bottomrule
\end{tabular}
\end{center}

\tpart{(4) Interpretation.} The coefficient is negative in every specification
and never approaches zero. Notably it \emph{strengthens} when eye care supply
and geography are added, because those covariates partly absorb the urban
clustering that otherwise masks it, and then attenuates modestly as social
vulnerability and screening completion enter. Its magnitude throughout is a
fraction of the social vulnerability coefficient estimated on the same sample.

The honest reading is that this is a small, stable, wrong-signed association
rather than a null. We do not interpret it as evidence that distance from
academic ophthalmology is protective, which is not clinically plausible. We
interpret it as the detection signature described in Tests~2 and~13: proximity
to subspecialty capacity raises the measured share of disease that is
vision-threatening by improving ascertainment, and no amount of covariate
adjustment removes a feature that is built into how the outcome is ascertained.

%=========================================================================
\section{Data provenance and reproducibility}
%=========================================================================

\tpart{(1) What is tested.} Whether the analytic file can be rebuilt from public
sources without manual steps, and whether derived quantities reproduce published
benchmarks.

\tpart{(2) How we tested it.} The pipeline is a numbered sequence of scripts that
downloads each source directly from its public endpoint, with no manual
download or point-and-click step. We validated the constructed outcome against
the published national estimates.

\tpart{(3) What we found.} Aggregating our county numerators reproduces the
published national figures: DR 26.43\% of adults with diabetes (published
26.43\%), VTDR 5.05\% (published 5.06\%), VTDR share 19.1\% (published
1.84/9.60 million $=$ 19.2\%). Source coverage was complete except as noted:
Puerto Rico is absent from VEHSS; 14 island and remote counties in Alaska and
Hawaii have no routable road path to any program.

\tpart{(4) Interpretation.} Two data-handling issues discovered during
construction are worth recording because they would silently corrupt a
replication. First, the earlier working dataset for this project used 2023
American Community Survey population, whose Connecticut planning-region
geography does not match the county geography of the 2021 surveillance and
vulnerability data; this produced eight apparently missing Connecticut counties
that are not missing at all once the correct vintage is used. Second, one
residency program (University of Wisconsin, Madison) geocoded to Madison County,
Illinois, roughly 350 miles away, which would have removed Wisconsin's flagship
program from the exposure surface for the upper Midwest. It was caught by
cross-checking program locations against Area Health Resources File ophthalmology
resident counts, a check we recommend to anyone geocoding program rosters.

\end{document}
